% ===== PRÉAMBULE CANONIQUE — à coller VERBATIM en tête de CHAQUE .tex =====
\documentclass[11pt,a4paper]{article}
\usepackage[utf8]{inputenc}\usepackage[T1]{fontenc}\usepackage{lmodern}
\usepackage[french]{babel}
\usepackage{amsmath,amssymb,mathtools}\usepackage{icomma}
\usepackage[version=4]{mhchem}          % \ce{...} : équations & formules
\usepackage{siunitx}\sisetup{locale=FR} % \qty{}{}, \si{}, \ang{}
\usepackage{chemfig}                    % \chemfig{...} : structures développées
\usepackage{xcolor}\usepackage{colortbl}
\usepackage{tikz}\usepackage{pgfplots}\pgfplotsset{compat=1.18}
\usepackage[most]{tcolorbox}\usepackage{enumitem}\usepackage{array,booktabs}
\usepackage[a4paper,margin=2.2cm]{geometry}
\usetikzlibrary{arrows.meta,calc,angles,quotes,positioning,patterns,backgrounds,shapes.geometric}
\definecolor{primary}{RGB}{222,49,99}\definecolor{primarydark}{RGB}{150,22,55}
\definecolor{defcol}{RGB}{0,114,178}\definecolor{methcol}{RGB}{52,92,148}
\definecolor{excol}{RGB}{0,135,90}\definecolor{attncol}{RGB}{200,40,55}
\tcbset{boxbase/.style={enhanced,breakable,boxrule=0.9pt,arc=2.5pt,
  left=3mm,right=3mm,top=2.3mm,bottom=2.3mm,fonttitle=\bfseries,coltitle=white}}
% --- boîtes TUTORIEL (noms EXACTS, ne pas renommer) ---
\newtcolorbox{cle}[1][]{boxbase,colback=defcol!6,colframe=defcol,
  title={\textsc{À retenir}\ifx\relax#1\relax\relax\else\textmd{ : \itshape #1}\fi}}
\newtcolorbox{methode}[1][]{boxbase,colback=methcol!7,colframe=methcol,sharp corners,
  title={\textsc{Méthode}\ifx\relax#1\relax\relax\else\textmd{ --- \itshape #1}\fi}}
\newtcolorbox{exemplebox}[1][]{boxbase,colback=excol!5,colframe=excol,
  title={\textsc{Exemple}\ifx\relax#1\relax\relax\else\textmd{ : \itshape #1}\fi}}
\newtcolorbox{piege}[1][]{boxbase,colback=attncol!6,colframe=attncol,
  title={\textsc{Piège}\ifx\relax#1\relax\relax\else\textmd{ : \itshape #1}\fi}}
\newtcolorbox{remarque}[1][]{boxbase,colback=black!4,colframe=black!45,
  title={\textsc{Remarque}\ifx\relax#1\relax\relax\else\textmd{ : \itshape #1}\fi}}
% --- boîtes EXERCICE / CORRIGÉ (noms EXACTS, auto-comptées) ---
\newtcolorbox[auto counter]{exo}[1][]{boxbase,colback=primary!4,colframe=primary,
  title={\textsc{Exercice}~\thetcbcounter\ifx\relax#1\relax\relax\else\textmd{ --- \itshape #1}\fi}}
\newtcolorbox[auto counter]{corrige}[1][]{boxbase,colback=excol!4,colframe=excol,
  title={\textsc{Corrigé}~\thetcbcounter\ifx\relax#1\relax\relax\else\textmd{ --- \itshape #1}\fi}}
\newcommand{\R}{\mathbb{R}}\newcommand{\N}{\mathbb{N}}\newcommand{\Z}{\mathbb{Z}}\newcommand{\ds}{\displaystyle}
% (un \usepackage{fancyhdr}+en-tête est possible ; il est ignoré côté web)
% =========================================================================
\begin{document}
\begin{center}\color{primarydark}\large\bfseries Thème 2 — Corrigé \quad$\bullet$\quad Niveau Facile\end{center}

\begin{corrige}[la partie électronique d'un couple]
$n=\mathrm{n.o.}(\text{Ox})-\mathrm{n.o.}(\text{Red})$ ; les électrons vont du côté de l'oxydant.
\begin{enumerate}[label=\alph*)]
  \item $n=3-2=1$ : \ce{Fe^{3+} + e^- <=> Fe^{2+}}.
  \item $n=2-0=2$ : \ce{Cu^{2+} + 2 e^- <=> Cu}.
  \item $n=2-0=2$ : \ce{Zn^{2+} + 2 e^- <=> Zn}.
  \item $n=3-0=3$ : \ce{Al^{3+} + 3 e^- <=> Al}.
\end{enumerate}
\end{corrige}

\begin{corrige}[demi-équations sans oxygène]
On équilibre les atomes, puis les électrons pour égaliser la charge.
\begin{enumerate}[label=\alph*)]
  \item \ce{Cl2 + 2 e^- <=> 2 Cl^-} (charge : gauche $-2$, droite $2\times(-1)=-2$).
  \item \ce{I2 + 2 e^- <=> 2 I^-}.
  \item \ce{Fe <=> Fe^{2+} + 2 e^-} (oxydation : électrons à droite).
  \item \ce{Ag+ + e^- <=> Ag}.
\end{enumerate}
\end{corrige}

\begin{corrige}[demi-équations avec oxygène, en milieu acide]
\begin{enumerate}[label=\alph*)]
  \item $n=5-2=3$. Il manque \textbf{2} O à droite $\Rightarrow$ $+2\,\ce{H2O}$ ; puis \textbf{4} H à
        gauche $\Rightarrow$ $+4\,\ce{H+}$ :
        \[\ce{NO3^- + 4 H+ + 3 e^- <=> NO + 2 H2O}\]
        Charge : gauche $-1+4-3=0$, droite $0$. \checkmark
  \item $n=6-4=2$. Il manque \textbf{2} O à droite $\Rightarrow$ $+2\,\ce{H2O}$ ; puis \textbf{4} H à
        gauche $\Rightarrow$ $+4\,\ce{H+}$ :
        \[\ce{SO4^{2-} + 4 H+ + 2 e^- <=> SO2 + 2 H2O}\]
        Charge : gauche $-2+4-2=0$, droite $0$. \checkmark
\end{enumerate}
\end{corrige}

\begin{corrige}[additionner deux demi-équations]
\begin{enumerate}[label=\alph*)]
  \item Les deux échangent 2 électrons : on additionne directement.
        \[\ce{Zn + Cu^{2+} <=> Zn^{2+} + Cu}\]
  \item La 2\textsuperscript{e} demi-équation échange 1 électron : on la multiplie par 2, puis on
        additionne :
        \[\ce{Cu + 2 Ag+ <=> Cu^{2+} + 2 Ag}\]
\end{enumerate}
\end{corrige}

\begin{corrige}[repérer les ions spectateurs]
\begin{enumerate}[label=\alph*)]
  \item L'ion \ce{SO4^{2-}} ne change pas (spectateur). Équation simplifiée :
        \ce{Zn + Cu^{2+} -> Zn^{2+} + Cu}.
  \item L'ion \ce{NO3^-} est spectateur. Équation simplifiée :
        \ce{Fe + 2 Ag+ -> Fe^{2+} + 2 Ag}.
\end{enumerate}
\end{corrige}

\end{document}
